% ==========================================================================
% FILE NAY DUOC SINH TU DONG -- KHONG SUA TAY.
% Sinh boi: code thuat toan/scripts/bao_cao/sinh_ket_qua_moi.py
% Sinh luc: 2026-09-21 10:40:30
% Muon doi noi dung: sua script roi chay lai.
% ==========================================================================

\begin{table}[H]
\centering\footnotesize
\caption{So sánh toàn bộ danh mục mô hình trên dữ liệu VN30F1M, ba phương pháp suy diễn}
\label{tab:buoc1-so-sanh}
\begin{tabular}{llrrrrrr}
\toprule
\textbf{Mô hình} & \textbf{PP} & $d$ & log-lik & AIC & BIC & QLIKE & NLS \\
\midrule
GARCH chuẩn & MLE & 4 & -2944,90 & 5897,80 & 5919,63 & 0,9836 & 1,6812 \\
 & VB-MF & 4 & -- & -- & -- & 0,9710 & 1,6680 \\
 & VB-FR & 4 & -- & -- & -- & 0,9852 & 1,6680 \\
GARCH Student-t & MLE & 5 & -2825,81 & 5661,62 & 5688,92 & 0,9943 & 1,5455 \\
 & VB-MF & 5 & -- & -- & -- & 1,0138 & 1,5461 \\
 & VB-FR & 5 & -- & -- & -- & 1,0082 & 1,5448 \\
GARCH-X Student-t & MLE & 7 & -2817,43 & 5648,87 & 5687,08 & 0,9520 & 1,5385 \\
 & VB-MF & 7 & -- & -- & -- & 1,5357 & 1,5529 \\
 & VB-FR & 7 & -- & -- & -- & 1,4223 & 1,5444 \\
EGARCH chuẩn & MLE & 5 & -2939,04 & 5888,09 & 5915,38 & 0,9625 & 1,6715 \\
 & VB-MF & 5 & -- & -- & -- & 0,9598 & 1,6624 \\
 & VB-FR & 5 & -- & -- & -- & 0,9651 & 1,6606 \\
EGARCH Student-t & MLE & 6 & -2822,04 & 5656,08 & 5688,83 & 0,9798 & 1,5376 \\
 & VB-MF & 6 & -- & -- & -- & 1,0005 & 1,5379 \\
 & VB-FR & 6 & -- & -- & -- & 0,9624 & 1,5371 \\
EGARCH-X Student-t & MLE & 8 & -2795,43 & 5606,87 & 5650,54 & 0,9234 & 1,5203 \\
 & VB-MF & 8 & -- & -- & -- & 0,9281 & 1,5215 \\
 & VB-FR & 8 & -- & -- & -- & 0,9258 & 1,5209 \\
MS-GARCH chuẩn & MLE & 10 & -2854,15 & 5728,29 & 5782,88 & 1,0232 & 1,5736 \\
 & VB-MF & 10 & -- & -- & -- & 1,0668 & 1,5721 \\
 & VB-FR & 10 & -- & -- & -- & 1,0539 & 1,5723 \\
MS-GARCH-X Student-t & MLE & 15 & -2810,29 & 5650,58 & 5732,46 & 0,9688 & 1,5373 \\
 & VB-MF & 15 & -- & -- & -- & 0,9978 & 1,5354 \\
 & VB-FR & 15 & -- & -- & -- & 1,1434 & 1,5392 \\
\bottomrule
\end{tabular}
\vspace{0.3em}
\raggedright\scriptsize
\emph{Ghi chú:} $d$ là số tham số tự do. Log-likelihood, AIC và BIC chỉ định nghĩa cho ước lượng điểm nên để trống ở hai phương pháp Bayes biến phân. QLIKE và NLS đo ngoài mẫu trên 435 phiên kiểm định; nhỏ hơn là tốt hơn. Mô hình chuẩn cài bằng chính khung Student-$t$ với $\nu$ cố định ở $10^6$, và $\nu$ đó không tính vào $d$.
\end{table}

\begin{table}[H]
\centering\footnotesize
\caption{Đóng góp riêng của từng thành phần cấu trúc, đo bằng chênh lệch AIC}
\label{tab:buoc1-phan-ra}
\begin{tabular}{lrrr}
\toprule
\textbf{Thành phần được thêm} & AIC trước & AIC sau & $\Delta$AIC \\
\midrule
Đuôi dày (chuẩn $\to$ Student-$t$), họ GARCH & 5897,80 & 5661,62 & -236,17 \\
Đuôi dày (chuẩn $\to$ Student-$t$), họ EGARCH & 5888,09 & 5656,08 & -232,01 \\
Biến ngoại sinh ($t \to$ X-$t$), họ GARCH & 5661,62 & 5648,87 & -12,75 \\
Biến ngoại sinh ($t \to$ X-$t$), họ EGARCH & 5656,08 & 5606,87 & -49,21 \\
Dạng hàm EGARCH, biến thể chuẩn & 5897,80 & 5888,09 & -9,71 \\
Dạng hàm EGARCH, biến thể Student-$t$ & 5661,62 & 5656,08 & -5,54 \\
Dạng hàm EGARCH, biến thể X Student-$t$ & 5648,87 & 5606,87 & -42,00 \\
Chuyển trạng thái, biến thể X Student-$t$ & 5648,87 & 5650,58 & 1,71 \\
\bottomrule
\end{tabular}
\vspace{0.3em}
\raggedright\scriptsize
\emph{Ghi chú:} $\Delta$AIC âm nghĩa là thành phần được thêm \emph{cải thiện} mô hình. Mọi giá trị lấy từ ước lượng hợp lý cực đại. Phép tách này chỉ hợp lệ vì danh mục được thiết kế thành một lưới hai chiều, trong đó mỗi cặp mô hình kề nhau chỉ khác \emph{đúng một} yếu tố.
\end{table}

\begin{table}[H]
\centering\footnotesize
\caption{Khôi phục tham số, MS-GARCH-X Student-$t$: độ phủ trung bình của khoảng tin cậy 95\%}
\label{tab:pra-coverage}
\begin{tabular}{llrrrr}
\toprule
\textbf{KB} & \textbf{Đặc điểm} & $n$ & MLE (Wald) & VB-MF & VB-FR \\
\midrule
$S_{1}$ & Hai trạng thái phân biệt rõ & 100 & 0,801 (50/100) & 0,323 & 0,469 \\
$S_{2}$ & Nhận dạng yếu, hai trạng thái gần nhau & 100 & 0,740 (49/100) & 0,495 & 0,567 \\
$S_{3}$ & Đuôi rất dày, $\nu_0 = 4$ & 100 & 0,792 (42/100) & 0,339 & 0,455 \\
\bottomrule
\end{tabular}
\vspace{0.3em}
\raggedright\scriptsize
\emph{Ghi chú:} Giá trị lý tưởng là $0{,}95$. Mỗi kịch bản $T = 2000$ quan sát, burn-in 200. Cột MLE là khoảng Wald từ Hessian quan sát, tính trên số lần lặp có Hessian xác định dương (trong ngoặc); ở khoảng một nửa số lần lặp nghiệm nằm trên vùng phẳng hoặc sống núi của hàm hợp lý nên Hessian không xác định dương và không có khoảng Wald. Nhãn trạng thái đã được sắp theo phương sai vô điều kiện $\omega_s/(1-\alpha_s-\beta_s)$ \emph{trước khi} gộp; không sắp nhãn thì mọi chỉ tiêu của tham số theo trạng thái đều vô nghĩa vì hai phương pháp Bayes đảo nhãn ở khoảng một nửa số lần lặp.
\end{table}

\begin{table}[H]
\centering\footnotesize
\caption{Khôi phục tham số, EGARCH-X Student-$t$ một trạng thái: hệ số đòn bẩy $\theta$}
\label{tab:prb-theta}
\begin{tabular}{llrrrr}
\toprule
\textbf{KB} & \textbf{Phương pháp} & $n$ & Độ chệch & RMSE & Độ phủ \\
\midrule
$E_{1}$ ($\theta_0 = -0{,}10$) & MLE & 100 & -0,0026 & 0,0185 & 0,950 \\
 & VB-MF & 100 & -0,0111 & 0,0229 & 0,890 \\
 & VB-FR & 100 & -0,0141 & 0,0259 & 0,950 \\
\addlinespace[0.2em]
$E_{2}$ ($\theta_0 = 0$) & MLE & 100 & -0,0006 & 0,0166 & 0,970 \\
 & VB-MF & 100 & -0,0001 & 0,0209 & 0,960 \\
 & VB-FR & 100 & -0,0026 & 0,0228 & 0,980 \\
\addlinespace[0.2em]
$E_{3}$ ($\nu_0 = 4$) & MLE & 100 & -0,0031 & 0,0216 & 0,950 \\
 & VB-MF & 100 & -0,0128 & 0,0253 & 0,910 \\
 & VB-FR & 100 & -0,0165 & 0,0286 & 0,960 \\
\addlinespace[0.2em]
\bottomrule
\end{tabular}
\vspace{0.3em}
\raggedright\scriptsize
\emph{Ghi chú:} $T = 2000$, $B = 100$ lần lặp mỗi ô, burn-in 200. Kịch bản $E_2$ kiểm định chiều \emph{ngược lại} của $E_1$: nếu $\hat\theta$ lệch khỏi 0 một cách hệ thống trên dữ liệu vốn không có đòn bẩy thì mọi kết luận về đòn bẩy của VN30F1M đều mất giá trị. Mô hình một trạng thái không có biến trạng thái ẩn nên điểm Brier không áp dụng.
\end{table}

\begin{table}[H]
\centering\footnotesize
\caption{EGARCH-X Student-$t$ trên VN30F1M: bốn phương pháp suy diễn trên cùng một hàm likelihood}
\label{tab:buoc3-bon-phuong-phap}
\begin{tabular}{lrrrrl}
\toprule
\textbf{Tham số} & MLE & VB-MF & VB-FR & MCMC & Khoảng tin cậy 95\% (MCMC) \\
\midrule
$\mu$ & 0,0796 & 0,0979 & 0,0919 & 0,0800 & [0,0336; 0,1226] \\
$\omega$ & 0,1599 & 0,1941 & 0,1897 & 0,1870 & [0,1093; 0,2948] \\
$\alpha$ & 0,0721 & 0,1081 & 0,1066 & 0,1121 & [0,0268; 0,2110] \\
$\theta$ & -0,1452 & -0,1407 & -0,1241 & -0,1399 & [-0,2125; -0,0710] \\
$\beta$ & 0,6832 & 0,6264 & 0,6390 & 0,6393 & [0,4912; 0,7562] \\
$\gamma_1$ & -0,0692 & -0,0877 & -0,0772 & -0,0765 & [-0,1268; -0,0370] \\
$\gamma_2$ & 0,2986 & 0,3228 & 0,3439 & 0,3324 & [0,2167; 0,4715] \\
$\nu$ & 4,2910 & 4,3897 & 4,3647 & 4,3644 & [3,5544; 5,4315] \\
\bottomrule
\end{tabular}
\vspace{0.3em}
\raggedright\scriptsize
\emph{Ghi chú:} MCMC dùng NUTS, 4 chuỗi $\times$ (1.000 warmup + 400 mẫu), khởi tạo phân tán quanh nghiệm MLE. Chẩn đoán: $\hat R$ lớn nhất 1,0125, ESS nhỏ nhất 1340, 0 phân kỳ, 8,1 giờ. Cột VB và MCMC là trung bình hậu nghiệm; cột MLE là ước lượng điểm.
\end{table}

\begin{table}[H]
\centering\footnotesize
\caption{Trả lời RQ3: tỷ số độ lệch chuẩn hậu nghiệm $R_{sd} = \mathrm{sd}_{\mathrm{VB}} / \mathrm{sd}_{\mathrm{MCMC}}$}
\label{tab:buoc3-rsd}
\begin{tabular}{lrrrrr}
\toprule
\textbf{Tham số} & sd MCMC & sd VB-MF & sd VB-FR & $R_{sd}$ VB-MF & $R_{sd}$ VB-FR \\
\midrule
$\mu$ & 0,0231 & 0,0219 & 0,0217 & 0,95 & 0,94 \\
$\omega$ & 0,0472 & 0,0175 & 0,0352 & 0,37 & 0,75 \\
$\alpha$ & 0,0463 & 0,0407 & 0,0470 & 0,88 & 1,01 \\
$\theta$ & 0,0362 & 0,0333 & 0,0361 & 0,92 & 1,00 \\
$\beta$ & 0,0691 & 0,0199 & 0,0516 & 0,29 & 0,75 \\
$\nu$ & 0,4809 & 0,3037 & 0,4652 & 0,63 & 0,97 \\
\bottomrule
\end{tabular}
\vspace{0.3em}
\raggedright\scriptsize
\emph{Ghi chú:} $R_{sd} < 1$ nghĩa là VB tự nhận chắc chắn hơn hậu nghiệm thật, lấy MCMC hội tụ làm chuẩn. Mean-field đánh giá thấp bất định của $\omega$ và $\beta$ khoảng ba lần --- đúng hai tham số nằm trên sống núi cong $\omega/(1-\beta)$ mà nghiên cứu mô phỏng đã chỉ ra; full-rank thu hẹp khoảng cách nhưng vẫn hẹp hơn 25\% ở đúng hai tham số đó.
\end{table}

